% !PRINTED 526
\weberpage{526}
%<<CONT
\gesperrt{so ist sie die symmetrische oder die alternirende Gruppe, und}

2. \gesperrt{wenn sie zwei Transpositionen mit einem gemeinsamen Element enthält, so ist sie die symmetrische Gruppe.}

Wenn nämlich in einer Gruppe der Cyklus $(0,1,2)$ vorkommt, und die Gruppe ausserdem transitiv ist, so enthält sie eine Permutation $\pi$, durch die 0 in 3 übergeführt wird. In dem dreigliedrigen Cyklus $\pi^{-1}(0,1,2)\pi$ kommt also gewiss die Ziffer 3 vor; er ist also etwa $(0,1,3)$, und demnach haben wir in der Gruppe
\[
\begin{gathered}
(0,1,2), \quad (0,1,2)^2 = (0,2,1) \\
(0,1,3), \quad (0,1,3)^2 = (0,3,1) \\
(0,1,2)(0,3,1) = (1,2,3) \\
(1,2,3)^2 = (1,3,2) \\
(0,3,1)(0,1,2) = (0,3,2) \\
(0,3,2)^2 = (0,2,3),
\end{gathered}
\]
d. h. alle überhaupt vorhandenen dreigliedrigen Cyklen. Demnach kommen alle Permutationen der alternirenden Gruppe in der fraglichen Gruppe vor (§. 153), und sie ist also entweder damit erschöpft, oder sie enthält noch eine Permutation der zweiten Art und ist dann die symmetrische Gruppe.

Enthält aber eine transitive Gruppe die beiden Transpositionen $(0,1)$, $(0,2)$, so enthält sie auch
\[
(0,1)(0,2) = (0,1,2),
\]
und damit die ganze alternirende Gruppe, und da sie auch eine Transposition enthält, so ist es die symmetrische Gruppe.

Wenn also nun eine von $P$ und $Q$ verschiedene Gruppe zunächst eine Transposition $(0,2)$ enthält, so muss sie, wenn sie transitiv ist, auch $(1,3)$ enthalten, und demnach die ganze intransitive Gruppe vierten Grades:
\weq{3}{1, \quad (0,2), \quad (1,3), \quad (0,2)(1,3).}
Soll also die Gruppe transitiv sein, so müssen noch weitere Permutationen dazu kommen. Diese können nur unter den viergliedrigen Cyklen oder unter den Transpositionspaaren gesucht werden. Ein viergliedriger Cyklus aber, in dem 0 und 2 cyklisch an einander stossen, kann nicht vorkommen, weil er, mit $(0,2)$ combinirt, einen dreigliedrigen Cyklus geben würde, wie:
\[
(0,2)(0,2,3,1) = (0,3,1).
\]
